\chapter{Progettazione del sistema}
In questo capitolo vengono illustrate l'architettura del sistema e le principali scelte progettuali adottate nello sviluppo del nuovo event display.
Tali scelte derivano dall’analisi dei requisiti discussa nel \cref{cap:software_esistente_requisiti} e mirano a garantire modularità, manutenibilità ed estendibilità del software.
Dopo una descrizione generale dell’interazione tra i moduli, verranno analizzate nel dettaglio le soluzioni progettuali adottate, con particolare attenzione ai \emph{design pattern} utilizzati \cite{Gamma_1994}.

% ====================================================
%
% 4.1 - ARCHITETTURA GENERALE
%
% ====================================================
\section{Architettura generale}
Per favorire la modularità del software, è stato scelto di dividerlo in package, così da ridurre le dipendenze che aumenterebbero la complessità della manutenzione.
La struttura generale è riportata in \cref{fig:diagramma_package}, che evidenzia come la gestione di ogni dipendenza esterna principale sia delegata a un solo modulo.

\begin{figure}[ht]
    \centering
    \includegraphics[width=\textwidth]{cap4/Package}
    \caption[Diagramma dei package del sistema]{Diagramma dei package del sistema, in cui sono riportate anche le principali dipendenze esterne}
    \label{fig:diagramma_package}
\end{figure}

Il cuore del sistema è contenuto nel package \texttt{core}, che racchiude le classi responsabili del ciclo di vita del software, dalla creazione della finestra all'interazione con l'utente.
Al suo interno è incluso il modulo \texttt{state}, a cui è delegata la gestione della macchina a stati finiti che governa il ciclo di vita dell'applicativo, che verrà illustrata nelle prossime sezioni.

Il modulo \texttt{io} isola la dipendenza da ROOT: contiene infatti le strutture utilizzate per rappresentare i dati degli eventi, e un manager che si interfaccia con sndsw per popolarle.
Permette quindi di concentrare in un unico punto la lettura e la conversione dei dati da ROOT, rendendone semplice il riadattamento in caso di future modifiche a sndsw.

Passando al livello di visualizzazione dei dati, il modulo \texttt{core} presenta una dipendenza verso \texttt{scene}, che modella la gerarchia tridimensionale (nodi, mesh, luci) e ne gestisce la rappresentazione.
Per evitare accoppiamenti rigidi, \texttt{scene} riceve i dati geometrici dal modulo \texttt{io} tramite strutture dati astratte, mentre definisce le proprietà cromatiche degli elementi attraverso il package \texttt{colors}.

La visualizzazione avviene grazie al modulo \texttt{rendering}, che è diviso in due sotto-moduli: \texttt{camera} che definisce come la scena viene inquadrata e proiettata e \texttt{engine}, che utilizza le API di OpenGL per eseguire il rendering degli elementi della scena.
Nell'eventualità di un upgrade futuro delle API grafiche, sarà quindi sufficiente riadattare le classi contenute in \texttt{engine}.
In realtà, anche il modulo \texttt{core} presenta una dipendenza limitata a OpenGL, circoscritta all'inizializzazione del contesto e alla gestione del render loop; tutta la logica di rendering degli oggetti è invece delegata ai moduli \texttt{rendering} e \texttt{scene}.

Infine, il package \texttt{ui} presenta il maggior numero di dipendenze, in quanto è responsabile delle funzionalità legate all’interazione con l’utente.
Per soddisfare i requisiti di usabilità, esso consente sia di visualizzare le informazioni testuali associate ai dati caricati, sia di controllare la rappresentazione della scena, ad esempio selezionando quali parti del rivelatore SND@LHC includere o escludere dalla visualizzazione.

% ====================================================
%
% 4.2 - GESTIONE DELLA MACCHINA A STATI FINITI
%
% ====================================================
\section{Gestione della macchina a stati finiti}
\label{sec:fsm}
L'intero software è basato su una macchina a stati finiti (\emph{Finite State Machine}, FSM) così da avere il pieno controllo sul suo flusso di esecuzione.
In questa sezione vengono presentati gli stati di alto livello, riportati in \cref{fig:diagramma_stati_generale}; successivamente saranno approfonditi gli stati più complessi, che verranno scomposti nei relativi sottostati.

\begin{figure}[H]
    \centering
    \includegraphics[width=0.85\textwidth]{cap4/Diagramma_FSM_generale}
    \caption[Diagramma della FSM del sistema]{Diagramma ad alto livello della FSM}
    \label{fig:diagramma_stati_generale}
\end{figure}

La prima divisione negli stati \texttt{RUNNING} e \texttt{CLOSED}, permette di chiudere il software in qualsiasi momento del suo ciclo di vita e centralizzarne l'arresto.
È ovvero possibile catturare l'evento di chiusura in svariati modi (chiusura finestra, pulsante nel menu, ecc.), garantendo in ogni caso la corretta deallocazione delle risorse e la corretta sequenza di arresto.

Il primo stato in cui si entra automaticamente all'avvio del sistema è quello di \texttt{INITIALIZATION}, in cui è richiesto all'utente di inserire numero di run ed evento per poter avviare il caricamento dei dati.
Eventuali errori vengono gestiti mediante lo stato \texttt{INIT\_ERROR}, che mostra all’utente la causa del problema e gli consente di riprendere l’esecuzione dal punto in cui si era interrotta.

Una volta completato il caricamento dei dati e inizializzati i componenti necessari, il sistema passa allo stato \texttt{INTERACTION}.
Qui avvengono le principali interazioni dell'utente con l'event display, tra cui il movimento della camera e il cambiamento delle opzioni di rendering.
È inoltre possibile esportare un’immagine della visualizzazione corrente tramite lo stato \texttt{EXPORT\_IMAGE}; al termine dell'esportazione, si ritorna automaticamente in interazione.

Durante l'interazione è possibile cambiare la run, l'evento o la geometria visualizzata nella scena: questo avviene grazie allo stato \texttt{CHANGE\_DATA}, che permette all'utente di caricare nuovi dati senza riavviare il software.
Come per l'inizializzazione, in caso di errori questi vengono gestiti nello stato \texttt{CHANGE\_ERROR}.

% ====================================================
% 4.2.1 - INIZIALIZZAZIONE
% ====================================================
\subsection{Inizializzazione}

\texttt{INITIALIZATION} è uno stato composto, il cui schema è riportato in \cref{fig:diagramma_stati_inizializzazione}.
Quando è necessario acquisire informazioni dall’utente, il processo viene suddiviso in due fasi: una fase di scelta (\texttt{CHOICE}), nella quale l’utente inserisce i dati richiesti, e una fase di caricamento (\texttt{LOAD}), nella quale il software elabora tali informazioni.
Questo schema viene adottato, in particolare, per l’inserimento della run, dell’evento e del file di geometria.

In ogni momento l’utente può tornare alla schermata di inserimento precedente tramite l’evento \texttt{goBack}, attivato mediante la pressione di un pulsante.

Per una data run, il sistema tenta inizialmente di individuare e caricare il file di geometria predefinito.
Qualora questa operazione dovesse fallire, ad esempio perché il file non esiste, viene mostrata all’utente una schermata di errore che consente di selezionare manualmente il file di geometria corretto.

Al termine dell'inizializzazione viene richiesto il caricamento della geometria al framework ROOT: questo avviene solo alla fine perché una volta caricata la geometria ROOT, non può più essere sostituita se non riavviando il software.
Se la geometria venisse caricata quando l'utente ha ancora la possibilità di cambiare run, nel caso in cui le due geometrie non dovessero corrispondere ROOT genererebbe un errore.

\begin{figure}[ht]
    \centering
    \includegraphics[width=0.97\textwidth]{cap4/Diagramma_FSM_init}
    \caption{Espansione dello stato composto \texttt{INITIALIZATION}}
    \label{fig:diagramma_stati_inizializzazione}
\end{figure}

In ogni caso, gli errori -- come visto nello schema generale di \cref{fig:diagramma_stati_generale} -- vengono gestiti tramite uno stato esterno detto \texttt{INIT\_ERROR}, in cui viene mostrata all'utente la causa scatenante (ad esempio run number inesistente), così che possa tornare indietro e rieseguire l'azione richiesta.


% ====================================================
% 4.2.2 - CAMBIO DATI
% ====================================================
\subsection{Cambio dati}

Per cambiare i dati visualizzati durante l'interazione si entra nello stato composto \texttt{CHANGE\_DATA} (il cui diagramma è riportato in \cref{fig:diagramma_stati_change_data}), che rispecchia l'inizializzazione: anche in questo caso, gli stati sono sdoppiati in \texttt{CHOICE} e \texttt{LOAD}, e gli errori sono gestiti entrando nello stato esterno \texttt{CHANGE\_ERROR}.
In questa fase si può però scegliere tramite interfaccia se cambiare il file di geometria, il numero di run o l'evento.
Poiché può essere necessario scorrere più eventi appartenenti alla stessa run, sarebbe poco efficiente richiedere ogni volta l’inserimento dell’intera coppia (run number, event number) e ricaricare dati già validi.
Per questo motivo, il sistema consente sia di specificare nuovamente entrambi i valori, sia di modificare soltanto il numero di evento nel caso in cui la run rimanga invariata.

\begin{figure}[H]
    \centering
    \includegraphics[width=0.88\textwidth]{cap4/Diagramma_FSM_change_data}
    \caption{Espansione dello stato composto \texttt{CHANGE\_DATA}}
    \label{fig:diagramma_stati_change_data}
\end{figure}

Il cambio della geometria è indipendente dal cambio di run ed evento: la collaborazione potrebbe infatti necessitare di caricare file di geometria 3D corrispondenti a configurazioni diverse del rivelatore, senza che questo implichi necessariamente un cambio dei dati fisici visualizzati.

% ====================================================
%
% 4.3 - STRUTTURA DELL'APPLICAZIONE
%
% ====================================================
\section{Struttura dell'applicazione}
Per garantire il corretto funzionamento del software è stata introdotta una classe principale, incaricata di coordinare i moduli fondamentali del sistema.
La classe \texttt{App}, posta all'interno del package \texttt{core}, rappresenta il punto di ingresso dell'applicazione e ne orchestra l'intero ciclo di vita, aggregando i componenti principali descritti in \cref{fig:classi_app}:
\begin{itemize}
    \item \texttt{AppStateManager}: mantiene lo stato corrente della FSM descritta nella \cref{sec:fsm} e ne gestisce i passaggi di stato.
    \item \texttt{AppSettings}: raccoglie tutte le impostazioni configurabili dall'utente, condivise tra i vari componenti.
    \item \texttt{WindowManager}: gestisce il ciclo di vita della finestra.
    \item \texttt{Callbacks}: raccoglie le callback per la gestione degli input da tastiera e mouse, delegandoli ai componenti appropriati.
    \item \texttt{Gui}: gestisce l'interfaccia utente.
    \item \texttt{Scene}: contiene la scena tridimensionale, con la geometria del rivelatore e le hit dell'evento corrente.
    \item \texttt{SndswEventManager}: si interfaccia con il framework sndsw per il caricamento dei dati fisici.
\end{itemize}

\begin{figure}[ht]
    \centering
    \includegraphics[width=\textwidth]{cap4/Classi_App}
    \caption{Dettaglio delle classi gestite da App}
    \label{fig:classi_app}
\end{figure}

Il metodo \texttt{run} contiene il render loop: cicla fino a che il software rimane in stato \texttt{RUNNING}, aggiornando tutti i componenti dell'app ed eseguendo il render di gui e scena tridimensionale.

Il metodo \texttt{update}, invocato ad ogni ciclo del render loop, funge da dispatcher della FSM: ad ogni frame legge lo stato corrente e delega l'operazione corrispondente al componente appropriato -- ad esempio il caricamento di una run a \texttt{SndswEventManager}.
Gli errori sollevati durante il caricamento vengono catturati e notificati all'\texttt{AppStateManager}, che provvede a transitare nello stato di errore appropriato.

% ====================================================
%
% 4.4 - DESIGN DETTAGLIATO
%
% ====================================================
\section{Design dettagliato}
In questa sezione vengono approfondite le principali scelte progettuali del sistema, illustrando per ciascuna il problema affrontato e la soluzione adottata.
La progettazione fa ampiamente uso dei \emph{design pattern}: per ogni componente viene descritto il pattern adottato e il contesto specifico della sua applicazione.

% ====================================================
% 4.4.1 - OGGETTI TRIDIMENSIONALI DELLA SCENA
% ====================================================
\subsection{Oggetti tridimensionali della scena}
\label{subsec:oggetti_tridimensionali_scena}
ROOT memorizza le geometrie dei rivelatori in formato proprietario.
Poiché la conversione diretta di tali file avrebbe comportato un lavoro considerevole, non strettamente legato agli obiettivi della tesi, si è scelto di adottare una conversione preventiva nel formato glTF 2.0 tramite uno strumento dedicato.
Questa scelta è giustificata anche dal fatto che le geometrie disponibili sono in numero limitato e vengono aggiornate raramente, indicativamente una volta all’anno.

\subsubsection{Lo standard glTF}

Il \emph{Graphics Library Transmission Format} (glTF) è uno standard open source per la condivisione di modelli e scene 3D, sotto forma di file di testo formattato in JSON.
Organizza internamente i dati secondo una struttura ad albero (in \cref{fig:gltf_spec} è riportato l’albero dello standard glTF 2.0), di cui sono rilevanti i seguenti nodi ai fini di questo progetto:
\begin{itemize}
    \item scene: contiene il puntatore al nodo radice dell’albero.
    \item node: generico nodo che contiene la matrice di trasformazione e opzionalmente una mesh e altri nodi figli.
    \item mesh: contenitore dei dati grezzi per il rendering, quali vertici, colori, normali e indici; eventualmente può anche avere un puntatore a un materiale.
    \item material: definisce le proprietà fisiche della superficie secondo il modello Physically Based Rendering.
\end{itemize}

\begin{figure}[ht]
    \centering
    \begin{subfigure}[b]{0.58\textwidth}
        \centering
        \includegraphics[width=\textwidth]{cap4/Struttura_glTF}
        \caption{Specifica glTF 2.0 \cite{Khronos_gltf_spec}}
        \label{fig:gltf_spec}
    \end{subfigure}
    \hfill
    \begin{subfigure}[b]{0.40\textwidth}
        \centering
        \includegraphics[width=\textwidth]{cap4/SceneObject}
        \caption{Oggetti della scena}
        \label{fig:diagramma_classi_object}
    \end{subfigure}
    \caption[Confronto tra specifica glTF e oggetti della scena]{Confronto tra la struttura dei file glTF e la rappresentazione degli oggetti della scena}
\end{figure}

Il \emph{Physically Based Rendering} (PBR) è un approccio che simula il comportamento fisico della luce basandosi su parametri quali il colore di base, il coefficiente metallico (metallic) e quello di rugosità (roughness).
Nel progetto, al contrario, il colore di base viene impiegato per creare un materiale -- classe \texttt{Material} -- secondo le componenti del modello di illuminazione di Blinn-Phong, discusso nella \cref{subsec:blinn_phong}.

\subsubsection{Struttura degli oggetti della scena}

Seguendo il pattern \emph{Composite}, sono state create le classi (illustrate in \cref{fig:diagramma_classi_object}) che riflettono la struttura ad albero glTF per la memorizzazione dei componenti della scena:
\begin{itemize}
    \item \texttt{Object} funge da componente astratto, ed è costituito dal solo nodo radice.
    \item \texttt{Node} rappresenta il componente composito, poiché può contenere altri nodi figli. Inoltre, mantiene il riferimento alle mesh associate al nodo e memorizza la matrice di modellazione corrispondente.
    \item \texttt{Mesh} è la foglia, in quanto rappresenta un elemento terminale della gerarchia privo di figli.
\end{itemize}

Poiché le geometrie possono essere condivise tra più nodi della scena, i dati delle mesh sono stati disaccoppiati dal concetto logico di mesh: un oggetto \texttt{Mesh} punta a una \texttt{GpuMesh} che può essere condivisa tra più istanze, evitando la duplicazione di dati geometrici in memoria GPU.
Lo stesso ragionamento vale per \texttt{Material}.
Questa soluzione corrisponde al pattern \emph{Flyweight}, applicato per minimizzare il consumo di memoria nella rappresentazione di geometrie e materiali ricorrenti.
Lo stato intrinseco è rappresentato dalle informazioni condivisibili (contenute in \texttt{GpuMesh} e \texttt{Material}), mentre lo stato estrinseco è dato dal contesto in cui l'oggetto viene utilizzato — ad esempio la matrice di modellazione contenuta in \texttt{Node}, che permette di posizionare e orientare la stessa \texttt{GpuMesh} in punti diversi della scena.
Per garantire la condivisione dello stato intrinseco, la creazione degli oggetti \texttt{Object} è delegata alla classe \texttt{ObjectFactory}, che svolge il ruolo di FlyweightFactory.
Per le geometrie di uso più frequente — ad esempio il cubo, utilizzato per rappresentare ogni singola hit del rivelatore — la \texttt{GpuMesh} viene istanziata una sola volta alla creazione della factory e successivamente condivisa tra tutti gli \texttt{Object} richiesti, evitando di duplicare i dati geometrici in memoria GPU anche quando il numero di istanze nella scena è elevato.

% ====================================================
% 4.4.2 - CARICAMENTO DEI DATI
% ====================================================
\subsection{Caricamento dei dati}
Per la lettura dei dati sperimentali tramite il framework ROOT è necessario gestire un elevato numero di classi.
In accordo con il pattern \emph{Facade}, il controllo di tutte queste dipendenze è stato centralizzato nella classe \texttt{SndswEventManager}, in modo da avere un unico punto di accesso verso l'esterno del software; questo disaccoppiamento risulta evidente nella \cref{fig:diagramma_classi_modulo_io}.

\begin{figure}[ht]
    \centering
    \includegraphics[width=\textwidth]{cap4/Modulo_io}
    \caption[Diagramma delle classi per la lettura dei dati]{Diagramma delle classi coinvolte nella lettura dei dati dal framework ROOT}
    \label{fig:diagramma_classi_modulo_io}
\end{figure}

\needspace{2\baselineskip}
Comunicando il numero di run prescelto ne vengono lette le caratteristiche e salvate nella classe \texttt{RunData}, così che l'event display possa mostrare all'utente i dati della run attualmente in memoria.

Il caricamento di un evento avviene invece passando le configurazioni come parametro tramite la struttura \texttt{ClusterConfiguration}, restituendo il risultato in un oggetto della classe \texttt{EventData}.
Come discusso nella \cref{sec:struttura_snd}, SND@LHC è composto da complessi sottorivelatori (classe \texttt{DetectorData}), ogni evento si compone delle hit (classe \texttt{HitData}) rilevate in ognuno di essi.

% ====================================================
% 4.4.3 - SCELTA TIPO DI PROIEZIONE
% ====================================================
\subsection{Scelta tipo di proiezione}
Per consentire all’utente di selezionare il tipo di proiezione da utilizzare, il sistema prevede due classi distinte, dedicate rispettivamente alla proiezione ortografica e alla proiezione prospettica.
Come mostrato nel capitolo precedente, entrambe seguono la stessa logica generale, ma differiscono nel modo in cui viene calcolata la matrice di proiezione.
Questa struttura ha motivato l’adozione del pattern \emph{Template Method}, che permette di definire uno schema comune delegando alle sottoclassi l’implementazione del calcolo specifico.
Gli attori del pattern sono illustrati in \cref{fig:diagramma_proiezioni}.

\begin{figure}[!ht]
    \centering
    \includegraphics[width=\textwidth]{cap4/Proiezioni}
    \caption[Diagramma delle classi per le proiezioni]{Diagramma delle classi per l'applicazione del pattern Template Method alle proiezioni}
    \label{fig:diagramma_proiezioni}
\end{figure}

Per migliorare le performance dell'event display evitando il ricalcolo continuo della matrice di proiezione, ne viene tenuta una copia sovrascritta ad ogni cambio delle caratteristiche della proiezione.
Poiché va ricalcolata solo al cambio dell'aspect ratio o del fov, i metodi template \texttt{setAspectRatio} e \texttt{setFov} definiti nella classe astratta \texttt{Basic\allowbreak{}Projection} richiamano il metodo astratto \texttt{computeProjectionMatrix}, ovvero l'operazione primitiva.
Le classi concrete, in questo caso \texttt{Ortho\allowbreak{}graphicProjec\allowbreak{}tion} e \texttt{Perspec\allowbreak{}tive\allowbreak{}Projection}, si occupano solo di specificare l'algoritmo di calcolo della matrice nel metodo astratto sovrascritto, evitando di ripetere codice comune e semplificando l'aggiunta di nuovi tipi di proiezioni.

Infine, nella proiezione ortografica è presente un metodo aggiuntivo per impostare la distanza della camera dal target: dato che questo tipo di proiezione preserva il parallelismo, il semplice avvicinamento della camera al target non produce l'effetto di zoom.
Si rende quindi necessario cambiare le dimensioni dei piani vicino e lontano del volume di vista.

% ====================================================
% 4.4.4 - COLORAZIONE HIT
% ====================================================
\subsection{Colorazione hit}
Per permettere all'utente di esplorare i dati da prospettive differenti, l'event display supporta diverse modalità di colorazione delle hit.

La determinazione del colore avviene sulla base di due criteri:
\begin{itemize}
    \item la grandezza fisica di riferimento, come tempo ed energia.
    \item la scala da adottare, ad esempio monocromatica, lineare o logaritmica.
\end{itemize}
La classe \texttt{ColorPalette} utilizza infatti due istanze del pattern \emph{Strategy} -- illustrato in \cref{fig:diagramma_colorpalette} -- in modo da disaccoppiare questa doppia variabilità senza violare il \emph{Single Responsibility Principle} (SRP), ma centralizzando l'operazione di colorazione delle hit.
Essa si occupa infatti di:
\begin{itemize}
    \item calcolare il colore per ogni hit.
    \item restituire il range di valori della scala di riferimento, comprensivo dell'unità di misura.
    \item generare campionamenti del colore su intervalli fissati, ad esempio per visualizzare la legenda della scala cromatica associata all’intervallo di valori considerato.
\end{itemize}

\begin{figure}[ht]
    \centering
    \includegraphics[width=\textwidth]{cap4/ColorStrategy}
    \caption[Diagramma delle classi per la colorazione delle hit]{Diagramma delle classi che illustra l'uso del doppio pattern Strategy per la gestione dei colori dell'event display}
    \label{fig:diagramma_colorpalette}
\end{figure}

\texttt{ColorPalette} svolge quindi il ruolo di classe contesto e utilizza due strategie di calcolo: \texttt{VariableGetter}, per selezionare la grandezza fisica di riferimento, e \texttt{ValueMapper}, per associare a una hit un colore in funzione del valore assunto all’interno di un determinato intervallo.

Nel sistema sono state implementate strategie concrete che consentono, tramite \texttt{VariableGetter}, di basare la colorazione sull’energia o sul tempo associati alla hit.
Tramite \texttt{ValueMapper}, invece, sono disponibili diverse modalità di costruzione della scala cromatica: monocromatica, lineare e logaritmica.

Questa struttura rende semplice l’estensione futura del sistema: nuovi criteri di colorazione potranno essere introdotti definendo ulteriori strategie concrete, senza modificare la logica interna di \texttt{ColorPalette}.

% ====================================================
% CONCLUSIONE CAPITOLO 4
% ====================================================
\bigskip
Le scelte progettuali descritte in questo capitolo costituiscono la base per l’implementazione del sistema.
Nel capitolo successivo verrà quindi illustrato come tali scelte siano state tradotte in una versione completa e funzionante del nuovo event display.